% Part: methods
% Chapter: proofs
% Section: proof-by-contradiction

\documentclass[../../../include/open-logic-section]{subfiles}

\begin{document}

\olfileid{mth}{prf}{con}

\olsection{বিরোধাভাসের সাহায্যে প্রমাণ}

প্রথমত, বিরোধাভাসের সাহায্যে প্রমাণ এমন একটি অনুমানের রীতি,
যা নঞর্থক দাবি প্রমাণে ব্যবহার করি। ধরো, দেখাতে চাও
কোনো দাবি~$p$ \emph{মিথ্যা}, অর্থাৎ~$\lnot p$। সবচেয়ে
ফলপ্রসূ পদ্ধতি হলো (a) $p$~সত্য ধরা, এবং (b) দেখানো
যে এই অনুমান থেকে এমন কিছু আসে, যা মিথ্যা বলে জানো।
``মিথ্যা বলে জানা কিছু'' হয়তো~$p$-র নিজের সঙ্গেই,
অথবা আলোচ্য সম্পূর্ণ দাবির অন্য কোনো পূর্বধারণার সঙ্গে
বিরোধ করে। যেমন, ``$q$ হলে $\lnot p$'' প্রমাণে
$q$~সত্য ধরে সেখান থেকে~$\lnot p$ প্রমাণ করতে
হয়। বিরোধাভাসের সাহায্যে $\lnot p$ প্রমাণ করলে
~$q$-র পাশাপাশি $p$-ও ধরতে হবে। $p$ থেকে
$\lnot q$ প্রমাণ করতে পারলে দেখিয়েছ, $p$
অনুমানটি অন্য অনুমান~$q$-র বিরোধী সিদ্ধান্তে নিয়ে
যায়; কারণ $q$ ও $\lnot q$ একই সঙ্গে সত্য হতে
পারে না। অবশ্য বিরোধে পৌঁছতে অন্য অনুমানের
রীতি ও সংজ্ঞা খোলাও ব্যবহার করতে হবে। উদাহরণ দেখি।

\begin{prop}
  $A \subseteq B$ এবং $B = \emptyset$ হলে $A$-র
  কোনো !!{element}s নেই।
\end{prop}

\begin{proof}
  $A \subseteq B$ এবং $B = \emptyset$ ধরি।
  দেখাতে চাই $A$-র কোনো !!{element}s নেই।
  \begin{quote}
    এটি শর্তাধীন দাবি বলে পূর্বপক্ষ ধরে উত্তরপক্ষ
    প্রমাণ করতে হবে। উত্তরপক্ষ: $A$-র কোনো
    !!{element}s নেই। আরও স্পষ্ট করে বললে,
    একটি~$x \in A$ আছে—এমন নয়।
  \end{quote}
  $A$-র কোনো !!{element}s নেই তখনই, যখন এমন
  কোনো~$x$ নেই যাতে $x \in A$।
  \begin{quote}
    তাই যা প্রমাণ করতে চাই, তা আসলে নঞর্থক
    দাবি~$\lnot p$: একটি $x \in A$ আছে—এমন
    নয়। বিরোধাভাসের সাহায্যে প্রমাণে সংশ্লিষ্ট
    ধনাত্মক দাবি~$p$, অর্থাৎ একটি $x \in A$ আছে,
    ধরে তা থেকে বিরোধ দেখাতে হবে। এই পদ্ধতি
    বোঝাতে ``বিরোধের জন্য ধরি'' বা শুধু ``উল্টোটা
    ধরি'' লিখে অনুমান~$p$ জানাই।
  \end{quote}
  উল্টোটা ধরি: একটি $x \in A$ আছে।
  \begin{quote}
    এবার নতুন এই অনুমান থেকে বিরোধ বের করব।
    আরও দুটি অনুমান আছে: $A \subseteq B$ এবং
    $B = \emptyset$। প্রথমটি থেকে $x \in B$ পাই:
  \end{quote}
  $A \subseteq B$ বলে $x \in B$।
  \begin{quote}
    কিন্তু $B = \emptyset$ বলে $B$-র প্রত্যেক
    !!{element} (যেমন $x$), $\emptyset$-রও
    !!a{element} হতে হবে।
  \end{quote}
  $B = \emptyset$ বলে $x \in \emptyset$। এটি
  বিরোধ, কারণ সংজ্ঞা অনুযায়ী $\emptyset$-র কোনো
  !!{element}s নেই।
  \begin{quote}
    এতেই প্রমাণ শেষ: সাজানো অনুমানগুলি থেকে
    কাঙ্ক্ষিত বিরোধ পেয়েছি; তাই সব অনুমান
    একসঙ্গে সত্য হতে পারে না। প্রথম দুটি অনুমান
    ($A \subseteq B$ এবং $B = \emptyset$) নিয়ে
    প্রশ্ন নেই। ফলে শেষে যোগ করা অনুমানটি—একটি
    $x \in A$ আছে—মিথ্যা। পূর্ণতার জন্য
    বিষয়টি স্পষ্ট করে লিখতে পারি।
  \end{quote}
  তাই একটি $x \in A$ আছে—এই অনুমান মিথ্যা;
  অতএব বিরোধাভাসের সাহায্যে প্রমাণে $A$-র
  কোনো !!{element}s নেই।
\end{proof}

প্রত্যেক ধনাত্মক দাবি একটি নঞর্থক দাবির সঙ্গে
সহজেই সমতুল্য: $p$ তখনই, যখন $\lnot\lnot p$।
তাই বিরোধাভাসের সাহায্যে ধনাত্মক দাবিও
``পরোক্ষভাবে'' প্রমাণ করা যায়। $p$ প্রমাণ করতে
তাকে $\lnot\lnot p$ হিসেবে পড়ো। $\lnot p$
থেকে বিরোধ দেখাতে পারলে বিরোধাভাসের সাহায্যে
$\lnot\lnot p$ এবং তাই~$p$ প্রতিষ্ঠিত হয়।

আগের উদাহরণে আমাদের লক্ষ্য ছিল নঞর্থক দাবি:
$A$-র কোনো !!{element}s নেই। তাই বিরোধের
জন্য নেওয়া অনুমানটি (একটি $x \in A$ আছে)
ধনাত্মক। এটি কাজে লাগার মতো একটি কল্পিত
$x \in A$ দিয়েছে; তাকে নিয়ে যুক্তি চালিয়ে
$x \in \emptyset$-এ পৌঁছেছি।

ধনাত্মক দাবি পরোক্ষভাবে প্রমাণ করলে বিরোধের জন্য
নেওয়া অনুমানটি নঞর্থক হবে। তবে প্রায়ই ধনাত্মক
দাবিকে নঞর্থক, আর নঞর্থক দাবিকে ধনাত্মক আকারে
সহজে বলা যায়। আগের প্রমাণটি প্রায় একই হতো,
যদি ``$A$-র কোনো !!{element}s নেই''-এর বদলে
``$A = \emptyset$'' প্রমাণ করতাম। ($=$-র সংজ্ঞায়
``$A = \emptyset$'' একটি সার্বিক দাবি: $A$-র
প্রত্যেক !!{element}, $\emptyset$-রও
!!{element}, এবং উল্টোটিও।) কিন্তু এটি
``একটি $x \in A$ আছে—এমন নয়'' নঞর্থক
দাবির সঙ্গে সহজেই সমতুল্য।

তাই কখনো~$p$ সরাসরি প্রমাণের চেয়ে অনুমান
হিসেবে~$\lnot p$ নিয়ে কাজ করা সহজ। সরাসরি
প্রমাণ সমান সহজ বা আরও সহজ হলেও (পরের
উদাহরণগুলির মতো), কেউ কেউ পরোক্ষ পথ পছন্দ করেন।
দ্বিনঞর্থকতা বিভ্রান্তিকর লাগলে মনে রেখো,
কোনো দাবির বিরোধাভাসের সাহায্যে প্রমাণ হলো তার
\emph{উল্টো} দাবি থেকে বিরোধের প্রমাণ। তাই
$\lnot p$-র বিরোধাভাসের সাহায্যে প্রমাণে~$p$
ধরে বিরোধ দেখাই; আর~$p$-র অনুরূপ প্রমাণে
~$\lnot p$ ধরে বিরোধ দেখাই।

\begin{prop}
$A \subseteq A \cup B$।
\end{prop}

\begin{proof}
  দেখাতে চাই $A \subseteq A \cup B$।
  \begin{quote}
    প্রথমে এটি ধনাত্মক দাবি: প্রত্যেক $x \in A$
    $A \cup B$-তেও আছে। তার নঞর্থকতা:
    কোনো $x \in A$ আছে, কিন্তু
    $\notin A \cup B$। তাই এই নঞর্থক দাবিটি
    ধরে দেখাব, তা থেকে বিরোধ আসে। এভাবেই
    পরোক্ষ প্রমাণ হবে।
  \end{quote}
  উল্টোটা ধরি, অর্থাৎ $A \nsubseteq A \cup B$।
  \begin{quote}
    $A \subseteq A \cup B$-র সংজ্ঞা আছে: প্রত্যেক
    $x \in A$, $\in A \cup B$-ও। এবার
    $A \nsubseteq A \cup B$-র অর্থ পেতে
    খোলা সংজ্ঞাটির উপর মৌলিক যৌক্তিক রূপান্তর
    করি: প্রত্যেক $x \in A$, $\in A \cup B$—
    কথাটি মিথ্যা তখনই, যখন \emph{কোনো}
    % BN-SRC-480: the negated membership concerns A union B, not C.
    $x \in A$-র জন্য $x \notin A \cup B$।
    (এখানে খুব সাবধান থাকতে হয়: ``সব $A$ হলো
    $B$'' ধরনের দাবির নঞর্থকতা ভুল বিশ্লেষণ
    করায় অনেক শিক্ষার্থীর বিরোধাভাসের প্রমাণ
    ব্যর্থ হয়।) অন্যভাবে,
    $A \nsubseteq A \cup B$ তখনই, যখন এমন
    একটি $x$ আছে যাতে $x \in A$ এবং
    $x \notin A \cup B$। এরপর সহজ।
  \end{quote}
  অতএব এমন $x \in A$ আছে যাতে
  $x \notin A \cup B$। $\cup$-র সংজ্ঞায়
  $x \in A \cup B$ তখনই, যখন $x \in A$
  অথবা $x \in B$। $x \in A$ বলে
  $x \in A \cup B$। কিন্তু অনুমানে ছিল
  $x \notin A \cup B$—এটি বিরোধ।
\end{proof}

\begin{prob}
$A \cap B \subseteq A$ পরোক্ষভাবে প্রমাণ করো।
\end{prob}

\begin{prop}
$A \subseteq B$ এবং $B \subseteq C$ হলে
$A \subseteq C$।
\end{prop}

\begin{proof}
  $A \subseteq B$ এবং $B \subseteq C$ ধরি।
  দেখাতে চাই $A \subseteq C$।
  \begin{quote}
    পরোক্ষ পথে যাই: যে কথাটি প্রতিষ্ঠা করতে চাই,
    তার নঞর্থকতা ধরি।
  \end{quote}
  উল্টোটা ধরি, অর্থাৎ $A \nsubseteq C$।
  \begin{quote}
    আগের মতো, $A \nsubseteq C$ তখনই, যখন
    প্রত্যেক $x \in A$ যে $\in C$-ও, তা
    সত্য নয়; অর্থাৎ কোনো $x \in A$,
    $\notin C$। অভ্যাস হলে এমন নঞর্থকতা
    খুলে বুঝতে আর বিশেষ ভাবতে হবে না।
  \end{quote}
  অন্যভাবে, একটি~$x$ আছে যাতে
  $x \in A$ এবং $x \notin C$।
  \begin{quote}
    এবার এটি ব্যবহার করে বিরোধে পৌঁছব।
    অন্য দুটি পূর্বধারণাও কাজে লাগবে।
  \end{quote}
  $A \subseteq B$ বলে $x \in B$।
  $B \subseteq C$ বলে $x \in C$।
  কিন্তু এটি $x \notin C$-র বিরোধী।
\end{proof}

\begin{prop}
$A \cup B = A \cap B$ হলে $A = B$।
\end{prop}

\begin{proof}
  $A \cup B = A \cap B$ ধরি। দেখাতে চাই
  $A = B$।
  \begin{quote}
    শুরুটা এখন পরিচিত:
  \end{quote}
  বিরোধের জন্য $A \neq B$ ধরি।
  \begin{quote}
    অনুমানটি $A \neq B$। $A = B$ তখনই,
    যখন $A \subseteq B$ এবং $B \subseteq A$;
    তাই $A \neq B$ তখনই, যখন
    $A \nsubseteq B$ \emph{অথবা}
    $B \nsubseteq A$। (নঞর্থকতা রূপান্তরে
    সতর্ক থাকা কত জরুরি, লক্ষ করো!{})
    এই বিযোজন থেকে বিরোধ দেখাতে ক্ষেত্রবিচার
    করে প্রতিটি ক্ষেত্রেই বিরোধ বের করব।
  \end{quote}
  $A \neq B$ তখনই, যখন $A \nsubseteq B$
  অথবা $B \nsubseteq A$। ক্ষেত্রবিচার করি।
  \begin{quote}
    প্রথম ক্ষেত্রে $A \nsubseteq B$ ধরি:
    কোনো $x$-এর জন্য $x \in A$ কিন্তু
    $\notin B$। $A \cap B$ হলো $A$ ও $B$-র
    সাধারণ !!{element}s-এর সেট; কোনো বস্তু
    এ দুটির একটিতে না থাকলে তাদের ছেদেও নেই।
    $A \cup B$-তে $A$ ও $B$-র সব সদস্য আছে;
    তাই যেকোনো একটিতে যা আছে, সংযোগেও তা
    আছে। ফলে $x \in A \cup B$ কিন্তু
    $x \notin A \cap B$; অতএব
    $A \cap B \neq A \cup B$।
  \end{quote}

  ক্ষেত্র 1: $A \nsubseteq B$। তাহলে কোনো
  $x$-এর জন্য $x \in A$ কিন্তু $x \notin B$।
  যেহেতু $x \notin B$, তাই $x \notin A \cap B$।
  $x \in A$ বলে $x \in A \cup B$। অতএব
  $A \cap B \neq A \cup B$; কিন্তু অনুমানে
  ছিল $A \cap B = A \cup B$। বিরোধ।

  ক্ষেত্র 2: $B \nsubseteq A$। তাহলে কোনো
  $y$-এর জন্য $y \in B$ কিন্তু $y \notin A$।
  আগের মতো $y \in A \cup B$, অথচ
  $y \notin A \cap B$। তাই আবার
  $A \cap B \neq A \cup B$; কিন্তু অনুমানে
  $A \cap B = A \cup B$। বিরোধ।
\end{proof}

\end{document}
